%!TEX root = ../../thesis.tex

\section{Ablations}\label{sec:results-ablations}

% Answers RQ3 (\cref{sec:introduction-rqs}). Every ablation is paired on the
% pre-registered set of RNG seeds (\cref{sec:methodology-runs}).

% TODO: Lead with the ablation that decides the thesis's central claim, then the rest.
%
% 1. Planning versus learned value: one MCTS simulation against the full budget.
%    If a single simulation nearly matches full search, the contribution is a learned
%    value function over coverage, NOT planning for fuzzing --- and that must be said
%    in those words. No claim about planning is made anywhere in this thesis without
%    this result. Run it on branches (\cref{sec:target-branches}) as well as sequence:
%    branches is the target built to need lookahead.
%
% 2. Reward design: absolute per-execution coverage against marginal new coverage
%    (\cref{sec:corpus-reward}); living reward only, against living plus terminal
%    coverage, against added novelty (\cref{sec:input-reward}). Prior work found the
%    history-based reward failed and the memoryless one worked; report whether that
%    replicates under a planner.
%
% 3. Observation design: cumulative and new coverage against each alone; action
%    history only; zero-based against raw.
%
% 4. Algorithm components: value prefix, the self-supervised consistency loss,
%    reanalyze fraction, prioritized replay. These test whether EfficientZero's
%    sample-efficiency machinery earns its keep in this domain rather than on Atari.
%
% 5. Search and capacity: simulation count, root noise, temperature, model size,
%    absorbing depth.
%
% 6. Corpus-specific factors (learned position, learned selection, learned retention)
%    are reported in \cref{sec:results-corpus-actions} where their context is.
%
% Close with a variance decomposition: how much outcome variance is attributable to
% the run, the hyperparameter and the target. "Run-to-run variance dominates" is a
% genuine finding,
% and isolating whether collapse tracks initialization or sampling distinguishes a
% bad-basin problem from an exploration problem (\cref{chap:failure-modes}).
